\chapter{دیدن جهان از دید پوسته}

در این فصل، به کالبدشکافی فرآیندهای پشت‌صحنه در رابط خط فرمان می‌پردازیم و رخدادهایی را بررسی می‌کنیم که پس از فشردن کلید \en{Enter} توسط پوسته \cmd{bash} مدیریت می‌شوند. علیرغم پیچیدگی مفاهیمی که با آن‌ها مواجه خواهیم شد، در این بخش تنها یک دستور جدید را معرفی می‌کنیم:

\begin{itemize}
  \item \cmd{echo}: نمایش یک سطر متنی در خروجی.
\end{itemize}

\section{مفهوم بسط (Expansion)}

هر زمان که دستوری را در خط فرمان وارد می‌کنید، پوسته \cmd{bash} پیش از آنکه فرآیند اجرا را آغاز کند، عملیاتی تحت عنوان بسط (\en{Expansion}) را روی متن ورودی اعمال می‌نماید. در طی این فرآیند، پوسته کاراکترهای رزرو شده و خاص را شناسایی کرده و آن‌ها را با مقادیر یا الگوهای متناظرشان جایگزین می‌کند. پیش‌تر با نمونه‌های اولیه‌ای از این سازوکار آشنا شده‌ایم؛ برای مثال، استفاده از نویسه جانشین (\en{Wildcard}) \cmd{*} که به مجموعه‌ای از نام‌ فایل‌ها بسط می‌یابد. به بیان ساده، در فرآیند بسط، شما عبارتی را تایپ می‌کنید و پوسته پیش از اقدام به اجرای دستور، آن را به عبارت دیگری دگرگون می‌کند.

برای تبیین دقیق‌تر مفهوم بسط، از دستور \cmd{echo} استفاده می‌کنیم که از فرامین داخلی (\en{Built-in}) پوسته محسوب می‌شود. این دستور عملکرد بسیار ساده‌ای دارد: هر آرگومان متنی که به آن ارجاع داده شود را مستقیماً در خروجی استاندارد چاپ می‌کند.

چنانچه یک رشته متنی ساده را به عنوان آرگومان به دستور \cmd{echo} ارجاع دهید، عیناً همان متن در خروجی چاپ می‌شود:

\begin{latin}
\begin{lstlisting}
[me@linuxbox ~]$ echo this is a test
this is a test
\end{lstlisting}
\end{latin}

اما پرسش اینجاست که در صورت استفاده از یک نویسه رزرو شده مانند \cmd{*} (\en{Asterisk})، پوسته چه واکنشی نشان می‌دهد؟

\begin{latin}
\begin{lstlisting}
[me@linuxbox ~]$ echo *
Desktop Documents ls-output.txt Music Pictures Public Templates Videos
\end{lstlisting}
\end{latin}

همان‌طور که ملاحظه می‌کنید، دستور \cmd{echo} به‌جای چاپ کاراکتر \cmd{*} در خروجی، فهرستی از نام تمامی فایل‌ها و دایرکتوری‌های موجود در مسیر جاری را نمایش داده است. این پدیده دقیقاً ناشی از فرآیند بسط (\en{Expansion}) توسط پوسته است. در واقع، پوسته پیش از آنکه کنترل را به دستور \cmd{echo} بسپارد، کاراکترهای خاصی نظیر \cmd{*} را شناسایی کرده و آن‌ها را با مقادیر واقعی موجود در سیستم فایل جایگزین می‌کند. به زبان فنی، دستور \cmd{echo} هرگز کاراکتر \cmd{*} را دریافت نمی‌کند، بلکه لیست کامل فایل‌ها را به عنوان آرگومان ورودی می‌پذیرد.

\section{بسط مسیر (Pathname Expansion)}

مکانیزمی که نویسه‌های جانشین (\en{Wildcards}) بر پایه آن عمل می‌کنند، بسط مسیر (\en{Pathname Expansion}) نامیده می‌شود. این فرآیند به پوسته اجازه می‌دهد تا الگوهای حاوی کاراکترهای خاص را به نام‌های واقعی و موجود در سیستم‌ فایل تبدیل کند. در حقیقت، بسیاری از تکنیک‌هایی که در فصول گذشته برای مدیریت و جست‌وجوی فایل‌ها به کار گرفتیم، بر پایه همین سازوکار بنا شده بودند.

فرض کنید ساختار دایرکتوری خانگی شما به شرح زیر است:

\begin{latin}
\begin{lstlisting}
[me@linuxbox ~]$ ls
Desktop    ls-output.txt  Pictures   Templates
Documents  Music          Public     Videos
\end{lstlisting}
\end{latin}

با استفاده از دستور \cmd{echo} می‌توانیم نحوه بسط الگوها توسط پوسته را مشاهده کنیم.

\textbf{جست‌وجوی فایل‌های آغازشده با حرف \en{D}:}

\begin{latin}
\begin{lstlisting}
[me@linuxbox ~]$ echo D*
Desktop Documents
\end{lstlisting}
\end{latin}

در اینجا، الگوی \cmd{D*} به نام دو دایرکتوری \file{Desktop} و \file{Documents} بسط یافته است.

\textbf{جست‌وجوی فایل‌های مختوم به حرف \en{s}:}

\begin{latin}
\begin{lstlisting}
[me@linuxbox ~]$ echo *s
Documents Pictures Templates Videos
\end{lstlisting}
\end{latin}

الگوی \cmd{*s} به تمامی اسامی که به نویسه \en{s} ختم می‌شوند، بسط می‌یابد.

\textbf{استفاده از کلاس‌های نویسه‌ای (\en{Character Classes}):}

\begin{latin}
\begin{lstlisting}
[me@linuxbox ~]$ echo [[:upper:]]*
Desktop Documents Music Pictures Public Templates Videos
\end{lstlisting}
\end{latin}

الگوی \cmd{[[:upper:]]*} به پوسته فرمان می‌دهد تا تمامی نام‌هایی را که با یک حرف بزرگ آغاز می‌شوند، شناسایی و بسط دهد.

\textbf{بسط الگو در مسیرهای سلسله‌مراتبی:}

\begin{latin}
\begin{lstlisting}
[me@linuxbox ~]$ echo /usr/*/share
/usr/kerberos/share /usr/local/share
\end{lstlisting}
\end{latin}

این الگو تمامی دایرکتوری‌های فرعی موجود در مسیر \file{/usr/} را که حاوی دایرکتوری‌ای به نام \file{share} هستند، یافته و نمایش می‌دهد. این مثال به‌وضوح نشان می‌دهد که بسط مسیر صرفاً به دایرکتوری جاری محدود نبوده و در تمامی سطوح سیستم فایل قابل اجرا است.

\section{بسط مسیر برای فایل‌های پنهان}

همان‌طور که می‌دانیم، در سیستم‌عامل لینوکس فایل‌هایی که نام آن‌ها با نقطه (\cmd{.}) آغاز می‌شود، به‌صورت پیش‌فرض به‌عنوان فایل‌های پنهان (\en{Hidden Files}) شناخته می‌شوند. مکانیسم بسط مسیر (\en{Pathname Expansion}) به‌طور پیش‌فرض از این قرارداد پیروی کرده و این فایل‌ها را در نتایج جست‌وجو لحاظ نمی‌کند. برای نمونه، دستور زیر فایل‌های پنهان را در خروجی خود نمایش نمی‌دهد:

\begin{latin}
\begin{lstlisting}
echo *
\end{lstlisting}
\end{latin}

برای بسط الگو به‌گونه‌ای که فایل‌های پنهان را نیز شامل شود، ممکن است در وهله اول شروع الگو با یک نقطه، منطقی به نظر برسد:

\begin{latin}
\begin{lstlisting}
echo .*
\end{lstlisting}
\end{latin}

این روش تا حدی هدف را محقق می‌کند، اما یک چالش جانبی به‌همراه دارد. با بررسی دقیق نتایج، متوجه خواهید شد که ارجاعات \cmd{.} (تک‌نقطه) و \cmd{..} (دونقطه) نیز در خروجی ظاهر می‌شوند که به‌ترتیب به «دایرکتوری جاری» و «دایرکتوری والد» اشاره دارند. حضور این دو مورد می‌تواند در فرآیندهای عملیاتی منجر به بروز خطاهای منطقی یا نتایج ناخواسته شود. برای مشاهده عینی این رفتار، می‌توانید دستور زیر را اجرا کنید:

\begin{latin}
\begin{lstlisting}
ls -d .* | less
\end{lstlisting}
\end{latin}

جهت بسط دقیق‌تر مسیر برای فایل‌های پنهان، می‌توان از یک الگوی هوشمندانه‌تر استفاده کرد:

\begin{latin}
\begin{lstlisting}
echo .[!.]*
\end{lstlisting}
\end{latin}

این الگو به پوسته (\en{Shell}) فرمان می‌دهد تا تنها نام‌هایی را بسط دهد که با یک نقطه آغاز می‌شوند، اما بلافاصله پس از آن، هر کاراکتری به‌جز نقطه قرار گرفته باشد. این روش برای تطبیق صحیح با فایل‌های پنهان (مانند فایل‌های پیکربندی) بسیار رایج است، هرچند همچنان فایل‌هایی را که نامشان با چندین نقطه شروع می‌شود، نادیده می‌گیرد.

به‌عنوان یک راهکار استاندارد و قابل‌اطمینان برای فهرست کردن فایل‌های پنهان، استفاده از دستور \cmd{ls} همراه با گزینه \cmd{-A} (مخفف \en{Almost all}) توصیه می‌شود. این دستور فهرست کاملی از تمامی محتویات، از جمله موارد پنهان را ارائه می‌دهد، در حالی که آگاهانه از نمایش ارجاعات \cmd{.} و \cmd{..} خودداری می‌کند:

\begin{latin}
\begin{lstlisting}
ls -A
\end{lstlisting}
\end{latin}

این رویکرد، بهینه‌ترین روش برای مدیریت دقیق و صحیح فایل‌های پنهان در محیط پوسته محسوب می‌شود.

\section{بسط تیلدا (Tilde Expansion)}

همان‌طور که در مبحث دستور \cmd{cd} اشاره شد، کاراکتر تیلدا (\cmd{\textasciitilde}) در محیط پوسته حامل معنای ویژه‌ای است. هنگامی که این نویسه در ابتدای یک عبارت قرار می‌گیرد، پوسته آن را به مسیر دایرکتوری خانگی (\en{Home Directory}) یک کاربر مشخص بسط می‌دهد. چنانچه پس از تیلدا نام کاربری درج نشود، پوسته به‌صورت پیش‌فرض آن را به مسیر دایرکتوری خانگی کاربر جاری بسط می‌دهد:

\begin{latin}
\begin{lstlisting}
[me@linuxbox ~]$ echo ~
/home/me
\end{lstlisting}
\end{latin}

در این نمونه، پوسته نویسه \cmd{\textasciitilde} را به آدرس مطلق دایرکتوری خانگی کاربر فعلی (در اینجا \en{me}) تبدیل کرده و دستور \cmd{echo} برآیند این پردازش را در خروجی چاپ می‌کند.

در صورتی که بلافاصله پس از تیلدا، نام کاربری خاصی قید شود، پوسته آن را به دایرکتوری خانگی همان کاربر اختصاصی نگاشت می‌کند. برای مثال، اگر کاربری با نام \en{bob} در سیستم تعریف شده باشد:

\begin{latin}
\begin{lstlisting}
[me@linuxbox ~]$ echo ~bob
/home/bob
\end{lstlisting}
\end{latin}

در این حالت، تیلدا به مسیر دایرکتوری خانگی کاربر \en{bob} بسط می‌یابد. این قابلیت، مکانیزمی سریع و استاندارد برای ارجاع به دایرکتوری‌های خانگی در سلسله‌مراتب سیستم فایل، بدون نیاز به دانستن مسیر دقیق آن‌ها، فراهم می‌آورد.

\section{بسط محاسباتی (Arithmetic Expansion)}

پوسته لینوکس این امکان را فراهم می‌آورد تا عملیات ریاضی را مستقیماً از طریق سازوکار بسط (\en{Expansion}) به انجام برسانید. این قابلیت، محیط خط فرمان را به یک ماشین‌حساب سریع و کاربردی تبدیل می‌کند.

ساختار استاندارد برای بسط محاسباتی به شرح زیر است:

\begin{latin}
\begin{lstlisting}
$((expression))
\end{lstlisting}
\end{latin}

در این نحو (\en{Syntax})، عبارت \file{expression} شامل اعداد و عملگرهای ریاضی است. پوسته ابتدا عبارت محصور در پرانتزهای دوگانه را محاسبه کرده و سپس حاصل نهایی را جایگزین کل عبارت اصلی می‌کند.

به مثال زیر توجه کنید:

\begin{latin}
\begin{lstlisting}
[me@linuxbox ~]$ echo $((2 + 2))
4
\end{lstlisting}
\end{latin}

در این نمونه، پوسته ابتدا عملیات \cmd{2 + 2} را پردازش کرده و سپس مقدار نهایی (\cmd{4}) را به‌عنوان آرگومان ورودی به دستور \cmd{echo} تحویل می‌دهد. این ویژگی برای انجام محاسبات سریع عددی در اسکریپت‌ها یا محیط تعاملی ترمینال، بدون نیاز به فراخوانی ابزارهای جانبی، بسیار کارآمد است.

بسط محاسباتی در پوسته صرفاً از اعداد صحیح (\en{Integers}) پشتیبانی می‌کند، اما با این وجود، طیف متنوعی از عملگرهای ریاضی را در اختیار کاربر قرار می‌دهد. در جدول زیر، عملگرهای اصلی و منطق عملکرد آن‌ها تشریح شده است:

\begin{table}[htbp]
\centering
\caption{عملگرهای محاسباتی (\en{Arithmetic Operators})}
\label{tab:arith-operators}
\begin{tabularx}{\textwidth}{c X}
\toprule
\textbf{عملگر} & \textbf{توضیح} \\
\midrule
\cmd{+} & جمع \\
\addlinespace
\cmd{-} & تفریق \\
\addlinespace
\cmd{*} & ضرب \\
\addlinespace
\cmd{/} & تقسیم (خروجی همواره به صورت عدد صحیح است) \\
\addlinespace
\cmd{\%} & باقی‌مانده تقسیم (\en{Modulo}) \\
\addlinespace
\cmd{**} & توان \\
\bottomrule
\end{tabularx}
\end{table}

در بسط محاسباتی، استفاده از فضاهای خالی (\en{Spaces}) اختیاری است و تاثیری در نتیجه نهایی ندارد. همچنین، پوسته از عبارت‌های تودرتو (\en{Nested Expressions}) نیز پشتیبانی می‌کند. برای مثال، جهت محاسبه‌ی حاصل ۵ به توان ۲ ضربدر ۳، می‌توان از ساختار زیر استفاده کرد:

\begin{latin}
\begin{lstlisting}
[me@linuxbox ~]$ echo $(($((5**2)) * 3))
75
\end{lstlisting}
\end{latin}

علاوه بر این، می‌توان همین نتیجه را با استفاده از یک جفت پرانتز دوگانه و مدیریت اولویت‌ها با پرانتزهای داخلی به دست آورد:

\begin{latin}
\begin{lstlisting}
[me@linuxbox ~]$ echo $(((5**2) * 3))
75
\end{lstlisting}
\end{latin}

در ادامه، مثالی از کاربرد عملگرهای تقسیم و باقی‌مانده (\en{Modulo}) ارائه شده است. به خاطر داشته باشید که در این محیط، تقسیم تنها نتایج عدد صحیح تولید می‌کند و بخش اعشاری نادیده گرفته می‌شود:

\begin{latin}
\begin{lstlisting}
[me@linuxbox ~]$ echo Five divided by two equals $((5/2))
Five divided by two equals 2
[me@linuxbox ~]$ echo with $((5%2)) left over.
with 1 left over.
\end{lstlisting}
\end{latin}

مبحث بسط محاسباتی و عملگرهای پیشرفته‌تر آن، در فصل ۳۴ به‌صورت جامع مورد بررسی قرار خواهد گرفت.

\section{بسط آکولاد (Brace Expansion)}

بسط آکولاد را می‌توان منعطف‌ترین و شاید خاص‌ترین نوع بسط در پوسته به‌شمار آورد. این قابلیت به شما اجازه می‌دهد تا از طریق یک الگوی فشرده حاوی آکولاد، چندین رشته متنی (\en{Strings}) مجزا تولید کنید.

ساختار فنی این بسط از سه جزء اصلی تشکیل شده است: یک پیشوند (\en{Preamble})، یک یا چند مجموعه محصور در آکولاد و یک پسوند (\en{postscript}). محتوای داخل آکولاد می‌تواند شامل یکی از موارد زیر باشد:

\begin{itemize}
  \item فهرستی از رشته‌ها که با ویرگول (\cmd{,}) از هم جدا می‌شوند.
  \item یک بازه (\en{Range}) از اعداد یا حروف که با دو نقطه (\cmd{..}) تعیین می‌شود.
\end{itemize}

\begin{note}
\textbf{نکته:} درج فاصله (\en{Space}) درون آکولادها مجاز نیست، مگر آنکه با استفاده از کوتیشن (\cmd{" "}) یا نویسه گریز (\en{Escape Character}) یعنی بک‌اسلش (\cmd{\textbackslash}) مهار شده باشد. نویسه گریز به پوسته اعلام می‌کند که نویسهٔ بعدی را نه به‌عنوان یک نویسهٔ خاص، بلکه به‌عنوان یک کاراکتر معمولی در نظر بگیرد. در غیر این صورت، پوسته فاصله را به‌عنوان پایان عبارت بسط تلقی کرده و فرآیند را نادرست انجام می‌دهد.
\end{note}

برای درک بهتر، ابتدا مثالی از تولید رشته‌های متنی ارائه می‌دهیم:

\begin{latin}
\begin{lstlisting}
[me@linuxbox ~]$ echo Front-{A,B,C}-Back
Front-A-Back Front-B-Back Front-C-Back
\end{lstlisting}
\end{latin}

در این سناریو، پوسته عبارت \cmd{\{A,B,C\}} را به سه مؤلفه مجزا بسط داده و هر یک را با پیشوند \cmd{Front-} و پسوند \cmd{-Back} ترکیب می‌کند.

این قابلیت در مدیریت فایل‌ها و ایجاد ساختارهای دایرکتوری پیچیده بسیار کارآمد است. برای مثال، جهت ایجاد هم‌زمان ۱۰ فایل تصویری با نام‌های \file{image01.jpg} تا \file{image10.jpg} می‌توان از دستور زیر بهره جست:

\begin{latin}
\begin{lstlisting}
[me@linuxbox ~]$ touch image{01..10}.jpg
\end{lstlisting}
\end{latin}

در این دستور، بازه عددی \cmd{\{01..10\}} به توالی اعداد از ۰۱ تا ۱۰ بسط یافته و سپس دستور \cmd{touch} برای ایجاد تک‌تک این فایل‌ها فراخوانی می‌شود. این ابزار از تکرار بیهوده دستورات جلوگیری کرده و بهره‌وری در خط فرمان را به شدت افزایش می‌دهد.

قابلیت بسط آکولاد (\en{Brace Expansion}) امکان تولید سریع توالی‌های عددی و نویسه‌ای را فراهم می‌کند. در نمونه زیر، نحوه ایجاد یک بازه عددی ساده را مشاهده می‌کنید:

\begin{latin}
\begin{lstlisting}
[me@linuxbox ~]$ echo Number_{1..5}
Number_1 Number_2 Number_3 Number_4 Number_5
\end{lstlisting}
\end{latin}

از نسخه ۴.۰ پوسته \en{Bash} به بعد، این ویژگی از صفرگذاری (\en{Zero-padding}) نیز پشتیبانی می‌کند که برای تولید نام فایل‌های هم‌تراز و مرتب‌شده در سیستم فایل بسیار حیاتی است:

\begin{latin}
\begin{lstlisting}
[me@linuxbox ~]$ echo {01..15}
01 02 03 04 05 06 07 08 09 10 11 12 13 14 15
[me@linuxbox ~]$ echo {001..15}
001 002 003 004 005 006 007 008 009 010 011 012 013 014 015
\end{lstlisting}
\end{latin}

این مکانیزم برای حروف الفبا نیز صادق است و حتی می‌توان توالی‌ها را به‌صورت معکوس تولید کرد:

\begin{latin}
\begin{lstlisting}
[me@linuxbox ~]$ echo {Z..A}
Z Y X W V U T S R Q P O N M L K J I H G F E D C B A
\end{lstlisting}
\end{latin}

علاوه بر این، بسط‌های آکولادی قابلیت تودرتو سازی (\en{Nesting}) دارند که امکان خلق الگوهای ترکیبی پیچیده را تنها با یک دستور واحد میسر می‌سازد:

\begin{latin}
\begin{lstlisting}
[me@linuxbox ~]$ echo a{A{1,2},B{3,4}}b
aA1b aA2b aB3b aB4b
\end{lstlisting}
\end{latin}

یکی از کاربردی‌ترین جنبه‌های این ابزار، خودکارسازی فرآیند ایجاد دایرکتوری‌ها (مدیریت انبوه دایرکتوری‌ها) است. برای مثال، جهت سازماندهی آرشیو تصاویر بر اساس سال و ماه، به‌جای درج دستی تک‌تک مسیرها که فرآیندی زمان‌بر و خطا‌پذیر است، می‌توان با یک دستور ساختار کامل را بنا نهاد. این رویکرد ضمن حفظ نظم زمانی، سرعت عملیاتی کاربر را در محیط پوسته به‌طور چشم‌گیری افزایش می‌دهد.

برای درک قدرت عملیاتی این قابلیت، سناریوی سازماندهی تصاویر را در نظر بگیرید. ابتدا یک دایرکتوری مرجع به نام \file{Photos} ایجاد کرده و به آن وارد می‌شویم. سپس با اجرای تنها یک دستور واحد، ساختار کاملی از دایرکتوری‌ها را برای بازه زمانی سال‌های ۲۰۰۷ تا ۲۰۰۹، به تفکیک ماه‌های ۰۱ تا ۱۲ بنا می‌کنیم:

\begin{latin}
\begin{lstlisting}
[me@linuxbox ~]$ mkdir Photos
[me@linuxbox ~]$ cd Photos
[me@linuxbox Photos]$ mkdir {2007..2009}-{01..12}
[me@linuxbox Photos]$ ls
2007-01  2007-07  2008-01  2008-07  2009-01  2009-07
2007-02  2007-08  2008-02  2008-08  2009-02  2009-08
2007-03  2007-09  2008-03  2008-09  2009-03  2009-09
2007-04  2007-10  2008-04  2008-10  2009-04  2009-10
2007-05  2007-11  2008-05  2008-11  2009-05  2009-11
2007-06  2007-12  2008-06  2008-12  2009-06  2009-12
\end{lstlisting}
\end{latin}

همان‌طور که مشاهده می‌شود، پوسته با پردازش یک عبارت کوتاه، به‌صورت خودکار ۳۶ دایرکتوری مجزا را با رعایت ترتیب زمانی و نام‌گذاری استاندارد ایجاد کرده است. این سطح از خودکارسازی، فرآیندهای اداری و مدیریتی در سیستم فایل را که پیش‌تر ملال‌آور و مستعد خطا بودند، به عملیاتی بسیار سریع و دقیق تبدیل می‌کند.

\section{بسط پارامتر (Parameter Expansion)}

در این بخش، تنها به‌صورت گذرا به مفهوم بسط پارامتر می‌پردازیم؛ چرا که کاربرد محوری این قابلیت بیشتر در حوزه اسکریپت‌نویسی پوسته مطرح است تا تعاملات مستقیم در خط فرمان. این ویژگی به شما اجازه می‌دهد قطعات کوچکی از داده‌ها را که تحت عنوان متغیر (\en{Variable}) شناخته می‌شوند، در حافظه ذخیره کرده و با تخصیص یک نام مشخص، در مراحل بعدی به آن‌ها ارجاع دهید.

برای نمونه، متغیری از پیش تعریف‌شده به‌نام \cmd{USER} در سیستم وجود دارد که نام کاربری شما را در خود نگه می‌دارد. جهت فراخوانی و مشاهده محتوای این متغیر، از نماد \cmd{\$} در ابتدای نام آن استفاده می‌کنیم:

\begin{latin}
\begin{lstlisting}
[me@linuxbox ~]$ echo $USER
me
\end{lstlisting}
\end{latin}

چنانچه تمایل دارید فهرست کاملی از متغیرهای تعریف‌شده در محیط عملیاتی (\en{Environment Variables}) خود را مشاهده کنید، می‌توانید از دستور \cmd{printenv} بهره بگیرید:

\begin{latin}
\begin{lstlisting}
[me@linuxbox ~]$ printenv | less
\end{lstlisting}
\end{latin}

نکته حائز اهمیت در فرآیند بسط پارامتر این است که در صورت تایپ نادرست نام یک متغیر، خروجی پوسته یک رشته خالی (\en{Null}) خواهد بود. این رفتار با سایر انواع بسط (مانند بسط مسیر) تفاوت دارد؛ چرا که در آن‌ها، عدم تطابق الگو معمولاً منجر به چاپ خودِ الگو می‌شود، اما در اینجا پوسته به‌سادگی از عبارت می‌گذرد:

\begin{latin}
\begin{lstlisting}
[me@linuxbox ~]$ echo $SUER

[me@linuxbox ~]$
\end{lstlisting}
\end{latin}

با اجرای دستور فوق، هیچ خروجی‌ای نمایش داده نمی‌شود؛ زیرا متغیری تحت عنوان \cmd{SUER} در سیستم تعریف نشده است. این ویژگیِ بسط پارامتر، روشی سریع برای احراز وجود یا عدم وجود یک متغیر در محیط پوسته به‌شمار می‌رود.

\section{جایگزینی فرمان (Command Substitution)}

مکانیزم جایگزینی فرمان به شما این امکان را می‌دهد که خروجی یک دستور را به‌عنوان یک واحد در فرآیند بسط (\en{Expansion}) به کار بگیرید. به بیان دقیق‌تر، پوسته ابتدا فرمان محصور در ساختار جایگزینی را اجرا نموده و سپس خروجی متنی حاصل را دقیقاً جایگزین عبارت اصلی می‌کند.

\begin{latin}
\begin{lstlisting}
[me@linuxbox ~]$ echo $(ls)
Desktop Documents ls-output.txt Music Pictures Public Templates Videos
\end{lstlisting}
\end{latin}

در این نمونه، پوسته ابتدا دستور \cmd{ls} را اجرا می‌کند؛ سپس خروجی آن را که فهرستی از نام فایل‌ها است، به‌عنوان آرگومان ورودی به دستور \cmd{echo} تحویل می‌دهد.

یکی از کاربردهای تخصصی این قابلیت زمانی است که قصد دارید مشخصات یک فایل اجرایی را بدون آگاهی از مسیر مطلق آن بررسی کنید. برای مثال، با بهره‌گیری از فرمان \cmd{which} می‌توان مسیر دقیق یک برنامه را استخراج کرد و آن را مستقیماً به دستور \cmd{ls -l} ارجاع داد:

\begin{latin}
\begin{lstlisting}
[me@linuxbox ~]$ ls -l $(which cp)
-rwxr-xr-x 1 root root 71516 2007-12-05 08:58 /bin/cp
\end{lstlisting}
\end{latin}

در این فرآیند، خروجی دستور \cmd{which cp} (یعنی مسیر \file{/bin/cp}) به‌عنوان آرگومان به \cmd{ls -l} منتقل شده و در نتیجه، جزئیات مالکیتی و دسترسی فایل \cmd{cp} نمایش داده می‌شود.

جایگزینی فرمان به دستورات ساده محدود نمی‌شود و می‌تواند کل یک پایپ‌لاین (\en{Pipeline}) را نیز شامل شود. این ویژگی زیربنای ایجاد ترکیب‌های عملیاتی بسیار قدرتمندی است:

\begin{latin}
\begin{lstlisting}
[me@linuxbox ~]$ file $(ls -d /usr/bin/* | grep zip)
/usr/bin/bunzip2:           symbolic link to `bzip2'
/usr/bin/bzip2:             ELF 32-bit LSB executable, Intel 80386, version 1 (SYSV), dynamically linked (uses shared libs), for GNU/Linux 2.6.9, stripped
/usr/bin/bzip2recover:      ELF 32-bit LSB executable, Intel 80386, version 1 (SYSV), dynamically linked (uses shared libs), for GNU/Linux 2.6.9, stripped
/usr/bin/funzip:            ELF 32-bit LSB executable, Intel 80386, version 1 (SYSV), dynamically linked (uses shared libs), for GNU/Linux 2.6.9, stripped
/usr/bin/gpg-zip:           Bourne shell script text executable
/usr/bin/gunzip:            symbolic link to `../../../bin/gunzip'
/usr/bin/gzip:              symbolic link to `../../../bin/gzip'
/usr/bin/mzip:              symbolic link to `mtools'
\end{lstlisting}
\end{latin}

در این سناریو، پوسته ابتدا پایپ‌لاین \cmd{ls -d /usr/bin/* | grep zip} را پردازش می‌کند. این زنجیره، فهرستی از فایل‌های موجود در دایرکتوری \file{/usr/bin/} را که نام آن‌ها حاوی عبارت \en{``zip''} است، تولید می‌نماید. در نهایت، خروجی این پردازش به‌عنوان آرگومان ورودی به دستور \cmd{file} داده می‌شود تا نوع و ساختار هر یک از آن فایل‌ها شناسایی و گزارش شود.

علاوه بر ساختار استاندارد، نحو دیگری برای جایگزینی فرمان وجود دارد که در پوسته‌های قدیمی‌تر (مانند \cmd{sh}) مرسوم بوده و همچنان در \cmd{bash} جهت حفظ سازگاری عقب‌رو (\en{Backward Compatibility}) پشتیبانی می‌شود. در این روش، به‌جای استفاده از ساختار \cmd{\$(...)} از علامت بک‌تیک یا بک‌کوت (\en{backquote}) استفاده می‌شود:

\begin{latin}
\begin{lstlisting}
[me@linuxbox ~]$ ls -l `which cp`
-rwxr-xr-x 1 root root 71516 2007-12-05 08:58 /bin/cp
\end{lstlisting}
\end{latin}

اگرچه این نحو هنوز معتبر و قابل اجرا است، اما در ادبیات مدرن لینوکس، استفاده از ساختار \cmd{\$(...)} اکیداً توصیه می‌شود. دلایل اصلی این ترجیح، خوانایی بالاتر و قابلیت تودرتوسازی (\en{Nesting}) آسان‌تر دستورات است؛ چرا که مدیریت چندین لایه از بک‌تیک‌ها در دستورات پیچیده بسیار دشوار و مستعد خطا است.

\section{سازوکار کوتیشن (Quoting)}

پس از آشنایی با انواع روش‌های بسط (\en{Expansion}) در پوسته، اکنون نوبت به یادگیری شیوه‌های کنترل و مهار این فرآیند می‌رسد. برای درک ضرورت این موضوع، به دو مثال زیر توجه کنید:

\begin{latin}
\begin{lstlisting}
[me@linuxbox ~]$ echo this is a     test
this is a test
\end{lstlisting}
\end{latin}

در مثال نخست، پوسته به‌صورت خودکار فواصل اضافی میان کلمات را حذف کرده و آن‌ها را به یک فاصله واحد تقلیل می‌دهد. این رفتار تحت عنوان تفکیک کلمات (\en{Word Splitting}) شناخته می‌شود.

\begin{latin}
\begin{lstlisting}
[me@linuxbox ~]$ echo The total is $100.00
The total is 00.00
\end{lstlisting}
\end{latin}

در مثال دوم، پوسته نویسه \cmd{\$} را به‌عنوان نشانه‌ای برای بسط پارامتر تفسیر می‌کند. از آنجا که متغیری تحت عنوان \cmd{1} در محیط تعریف نشده است، پوسته آن را با یک رشته تهی جایگزین کرده و تنها باقی‌مانده عبارت را چاپ می‌کند.

برای پیش‌گیری از این رفتارهای ناخواسته، پوسته از مکانیزمی به نام کوتیشن (\en{Quoting}) بهره می‌برد. این قابلیت به کاربر اجازه می‌دهد تا به‌صورت گزینشی، فرآیندهای بسط را متوقف کرده و به پوسته فرمان دهد که برخی نویسه‌ها یا رشته‌های متنی را دقیقاً به‌صورت تحت‌اللفظی (\en{Literal}) و بدون تفسیر معنایی اجرا کند.

\subsection{دابل‌کوتیشن (Double Quotes)}

یکی از پرکاربردترین روش‌های مهار بسط در پوسته، استفاده از دابل‌کوتیشن (\en{Double Quotes}) است. هنگامی که متنی درون دابل‌کوتیشن (\cmd{" "}) محصور می‌شود، اکثر نویسه‌های خاصِ پوسته معنای عملیاتی خود را از دست داده و به‌عنوان کاراکترهای متنیِ ساده تفسیر می‌گردند. این سازوکار فرآیندهایی نظیر تفکیک کلمات، بسط مسیر، بسط تیلدا و بسط آکولاد را به‌طور کامل متوقف می‌کند.

با این وجود، دابل‌کوتیشن یک مهارکننده مطلق نیست؛ نویسه‌های \cmd{\$}، \cmd{\textbackslash} و \cmd{`} (بک‌تیک) همچنان قابلیت عملیاتی خود را حفظ می‌کنند. به این معنا که بسط پارامتر، بسط محاسباتی و جایگزینی فرمان حتی درون دابل‌کوتیشن نیز فعال باقی می‌مانند.

یکی از حیاتی‌ترین کاربردهای دابل‌کوتیشن، مدیریت فایل‌هایی است که در نام خود دارای فاصله (\en{Space}) هستند. به‌صورت پیش‌فرض، پوسته هر فاصله را به‌عنوان جداکننده آرگومان‌ها در نظر می‌گیرد. برای مثال، اگر فایلی تحت عنوان \file{two words.txt} داشته باشید، اجرای دستور زیر با خطا مواجه می‌شود؛ چرا که پوسته تصور می‌کند شما به‌دنبال دو فایل مجزا به نام‌های \file{two} و \file{words.txt} هستید:

\begin{latin}
\begin{lstlisting}
[me@linuxbox ~]$ ls -l two words.txt
ls: cannot access two: No such file or directory
ls: cannot access words.txt: No such file or directory
\end{lstlisting}
\end{latin}

برای اصلاح این رفتار، باید نام فایل را درون دابل‌کوتیشن قرار دهید تا پوسته کل عبارت را به‌عنوان یک واحد آرگومانی منفرد پردازش کند:

\begin{latin}
\begin{lstlisting}
[me@linuxbox ~]$ ls -l "two words.txt"
-rw-rw-r-- 1 me me 18 2016-02-20 13:03 two words.txt
\end{lstlisting}
\end{latin}

با بهره‌گیری از این تکنیک، می‌توانید فایل مذکور را اصلاح کرده و از چالش‌های نام‌گذاری غیراستاندارد رها شوید:

\begin{latin}
\begin{lstlisting}
[me@linuxbox ~]$ mv "two words.txt" two_words.txt
\end{lstlisting}
\end{latin}

با جایگزینی فاصله با نویسه زیرخط (\cmd{\_})، دیگر نیازی به استفاده مداوم از کوتیشن برای فراخوانی این فایل نخواهید داشت.

به خاطر داشته باشید که در محیط دابل‌کوتیشن (\en{Double Quotes})، فرآیندهای بسط پارامتر، بسط محاسباتی و جایگزینی فرمان همچنان فعال هستند و طبق منطق خود عمل می‌کنند:

\begin{latin}
\begin{lstlisting}
[me@linuxbox ~]$ echo "$USER $((2+2)) $(df -h)"
me 4
Filesystem      Size  Used Avail  Use% Mounted on
tmpfs           1.6G  2.0M  1.6G    1% /run
/dev/sda2        94G   19G   71G   21% /
tmpfs           7.8G     0  7.8G    0% /dev/shm
tmpfs           5.0M  4.0K  5.0M    1% /run/lock
/dev/sda1       975M  6.1M  969M    1% /boot/efi
/dev/sdb1       907G  574G  287G   67% /home
tmpfs           1.6G  1.8M  1.6G    1% /run/user/1000
\end{lstlisting}
\end{latin}

در این مثال، پوسته ابتدا سه فرآیند بسط مجزا را به موازات هم پردازش می‌کند:

\begin{itemize}
  \item \textbf{بسط پارامتر:} عبارت \cmd{\$USER} به نام کاربری (در اینجا \en{me}) بسط می‌یابد.
  \item \textbf{بسط محاسباتی:} عبارت \cmd{\$((2+2))} به عدد \cmd{4} ترجمه می‌شود.
  \item \textbf{جایگزینی فرمان:} عبارت \cmd{\$(df -h)} با خروجی متنیِ حاصل از اجرای دستور \cmd{df -h} جایگزین می‌گردد.
\end{itemize}

در نهایت، تمامی این نتایج به‌دلیل محصور بودن در دابل‌کوتیشن، به‌عنوان یک رشته متنی واحد (بدون در نظر گرفتن فواصل داخلی به‌عنوان جداکننده آرگومان) به دستور \cmd{echo} تحویل داده می‌شوند.

همان‌طور که پیش‌تر ملاحظه شد، فرآیند تفکیک کلمات (\en{Word Splitting}) به‌صورت پیش‌فرض منجر به حذف فواصل اضافی در خط فرمان می‌گردد:

\begin{latin}
\begin{lstlisting}
[me@linuxbox ~]$ echo this is a     test
this is a test
\end{lstlisting}
\end{latin}

در تنظیمات استاندارد، پوسته (\en{Shell}) هرگونه فاصله (\en{Space})، تب (\en{Tab}) یا کاراکتر خط جدید (\en{Newline}) را به‌عنوان جداکننده (\en{Delimiter}) میان کلمات تلقی می‌کند. در مثال فوق، پوسته کلمات \file{this}، \file{is}، \file{a} و \file{test} را به‌عنوان چهار آرگومان مجزا تشخیص داده و به دستور \cmd{echo} ارسال می‌کند. از آنجا که فواصل میان این کلمات صرفاً نقش تفکیک‌کننده آرگومان‌ها را ایفا کرده‌اند، در خروجی نهایی حذف شده و با یک فاصله واحد جایگزین می‌شوند.

با این حال، با بهره‌گیری از دابل‌کوتیشن (\en{Double Quotes})، می‌توان این رفتار پیش‌فرض را مهار کرد:

\begin{latin}
\begin{lstlisting}
[me@linuxbox ~]$ echo "this is a     test"
this is a     test
\end{lstlisting}
\end{latin}

در این وضعیت، به‌دلیل محصور شدن متن در دابل‌کوتیشن، پوسته تمامی فواصل را به‌عنوان بخشی لاینفک از خودِ آرگومان در نظر گرفته و از حذف آن‌ها خودداری می‌کند. در نتیجه، خط فرمان تنها شامل یک دستور و یک آرگومان واحد است که دقیقاً با همان ساختار متنیِ ورودی مطابقت دارد.

واقعیتِ تلقی شدنِ نویسه «خطِ جدید» (\en{Newline}) به‌عنوان جداکننده در فرآیند تفکیک کلمات (\en{Word Splitting})، تأثیر عمیق و ملموسی بر خروجیِ حاصل از جایگزینی فرمان (\en{Command Substitution}) دارد. برای درک این تمایز فنی، نمونه زیر را بررسی می‌کنیم:

\textbf{سناریوی اول: اجرای دستور بدون استفاده از کوتیشن}

\begin{latin}
\begin{lstlisting}
[me@linuxbox ~]$ echo $(df -h)
Filesystem Size Used Avail Use% Mounted on tmpfs 1.6G 2.0M 1.6G 1% /run /dev/sda2 94G 19G 71G 21% / tmpfs 7.8G 0 7.8G 0% /dev/shm tmpfs 5.0M 4.0K 5.0M 1% /run/lock /dev/sda1 975M 6.1M 969M 1% /boot/efi /dev/sdb1 907G 574G 287G 67% /home tmpfs 1.6G 1.8M 1.6G 1% /run/user/1000
\end{lstlisting}
\end{latin}

در این وضعیت، پوسته ابتدا دستور \cmd{df -h} را فراخوانی می‌کند. سپس خروجی چندخطی و ساختاریافته آن را دریافت کرده و تمامی فواصل (\en{Spaces}) و خطوط جدید (\en{Newlines}) موجود در آن را به‌عنوان جداکننده (\en{Delimiter}) برای دستور \cmd{echo} تفسیر می‌کند.

در نتیجه، ساختار ستونی خروجی از هم پاشیده و به یک رشته متنیِ پیوسته و تک‌خطی تبدیل می‌شود. در این فرآیند، هر کلمه به یک آرگومان مستقل بدل می‌گردد؛ به‌طوری که در مثال فوق، خروجی به ۴۹ آرگومان مجزا شکسته شده و به دستور \cmd{echo} تحویل داده می‌شود.

\textbf{سناریوی دوم: استفاده از دابل‌کوتیشن (\en{Double Quotes})}

\begin{latin}
\begin{lstlisting}
[me@linuxbox ~]$ echo "$(df -h)"
Filesystem      Size  Used Avail Use% Mounted on
tmpfs          1.6G  2.0M  1.6G   1% /run
/dev/sda2       94G   19G   71G  21% /
tmpfs          7.8G     0  7.8G   0% /dev/shm
tmpfs          5.0M  4.0K  5.0M   1% /run/lock
/dev/sda1      975M  6.1M  969M   1% /boot/efi
/dev/sdb1      907G  574G  287G  67% /home
tmpfs          1.6G  1.8M  1.6G   1% /run/user/1000
\end{lstlisting}
\end{latin}

در این حالت، به‌کارگیری دابل‌کوتیشن منجر به غیرفعال شدن مکانیسم تفکیک کلمات (\en{Word Splitting}) می‌گردد. در نتیجه، پوسته کل خروجی دستور \cmd{df -h} را به‌عنوان یک رشته متنی واحد (\en{Single String}) تلقی می‌کند که تمامی فواصل و نویسه‌های خط جدید (\en{Newline}) اصلی خود را حفظ کرده است.

بنابراین، دستور \cmd{echo} تمام محتوا را با ساختار و فرمت‌بندی اولیه چاپ می‌کند؛ گویی تنها یک آرگومان منفرد به آن ارسال شده است. این نمونه به‌خوبی نشان می‌دهد که استفاده از دابل‌کوتیشن تا چه حد برای حفظ یکپارچگی و خوانایی خروجیِ یک فرمان در محیط پوسته حیاتی است.

\subsection{تک‌کوتیشن (Single Quotes)}

جهت ابطال و سرکوب مطلق تمامی انواع بسط (\en{Expansion})، از تک‌کوتیشن (\cmd{' '}) استفاده می‌شود. هر عبارتی که درون این تک‌کوتیشن قرار گیرد، فاقد هرگونه ارزش عملیاتی برای پوسته تلقی شده و صرفاً به‌صورت تحت‌اللفظی (\en{Literal}) تفسیر می‌گردد؛ به‌طوری که پوسته هیچ‌گونه پردازشی روی آن انجام نمی‌دهد.

برای درک عمیق تمایز میان سطوح مختلف کوتیشن، سه وضعیت زیر را با یکدیگر مقایسه می‌کنیم:

\begin{latin}
\begin{lstlisting}
[me@linuxbox ~]$ echo text ~/*.txt {a,b} $(echo foo) $((2+2)) $USER
text /home/me/ls-output.txt a b foo 4 me
[me@linuxbox ~]$ echo "text ~/*.txt {a,b} $(echo foo) $((2+2)) $USER"
text ~/*.txt {a,b} foo 4 me
[me@linuxbox ~]$ echo 'text ~/*.txt {a,b} $(echo foo) $((2+2)) $USER'
text ~/*.txt {a,b} $(echo foo) $((2+2)) $USER
\end{lstlisting}
\end{latin}

تحلیل رفتار پوسته در هر لایه:

\begin{itemize}
  \item \textbf{بدون کوتیشن:} تمامی بسط‌ها (تیلدا، مسیر، آکولاد، جایگزینی فرمان، محاسباتی و پارامتر) به‌طور کامل اجرا می‌شوند.
  \item \textbf{دابل‌کوتیشن (\en{Double Quotes}):} بسط‌های جایگزینی فرمان، محاسباتی و پارامتر همچنان فعال هستند، اما بسط‌های تیلدا، مسیر فایل و آکولاد مهار می‌شوند.
  \item \textbf{تک‌کوتیشن (\en{Single Quotes}):} تمامی قابلیت‌های بسط به‌طور کامل متوقف شده و متن دقیقاً با همان ساختار نوشتاری در خروجی چاپ می‌شود.
\end{itemize}

همان‌طور که ملاحظه می‌شود، با افزایش سطح سخت‌گیری در کوتیشن، تعداد بیشتری از ویژگی‌های تفسیریِ پوسته نادیده گرفته می‌شوند تا دسترسی به محتوای خام و اصلی میسر گردد. این سلسله‌مراتب، کنترل دقیقی بر نحوه پردازش دستورات را برای کاربر فراهم می‌آورد.

\section{نویسه گریز (Escape Character)}

گاهی اوقات به‌جای محصور کردن کل یک رشته در کوتیشن، تنها نیاز داریم معنای عملیاتی یک کاراکتر خاص را خنثی کنیم. برای این منظور از نویسه گریز (\en{Escape Character}) یعنی بک‌اسلش (\cmd{\textbackslash}) استفاده می‌کنیم. این نویسه به پوسته فرمان می‌دهد که کاراکتر بلافاصله پس از خود را به‌صورت تحت‌اللفظی (\en{Literal}) و بدون در نظر گرفتن نقش فنی آن تفسیر کند.

این روش به‌ویژه در محیط دابل‌کوتیشن (\en{Double Quotes}) برای مهار انتخابیِ برخی بسط‌ها بسیار پرکاربرد است. به مثال زیر توجه کنید:

\begin{latin}
\begin{lstlisting}
[me@linuxbox ~]$ echo "The balance for user $USER is: \$5.00"
The balance for user me is: $5.00
\end{lstlisting}
\end{latin}

در این نمونه، استفاده از \cmd{\textbackslash\$} به پوسته ابلاغ می‌کند که علامت دلار دوم را صرفاً به‌عنوان یک نماد متنی چاپ کند. این در حالی است که عبارت \cmd{\$USER} به‌دلیل عدم استفاده از بک‌اسلش، همچنان طبق روال معمول به نام کاربری (در اینجا \en{me}) بسط می‌یابد.

استفاده از نویسه گریز (\en{Escape Character}) یکی از روش‌های رایج برای درج کاراکترهای ویژه در نام فایل‌ها است. در محیط پوسته، نویسه‌هایی نظیر \cmd{\$}، \cmd{!}، \cmd{\&} و حتی «فاصله»، حامل معانی عملیاتی خاصی هستند. برای اینکه پوسته این کاراکترها را نه به‌عنوان یک فرمان، بلکه به‌عنوان بخشی از نام فایل تلقی کند، باید پیش از هر یک از آن‌ها یک بک‌اسلش (\cmd{\textbackslash}) درج کرد:

\begin{latin}
\begin{lstlisting}
[me@linuxbox ~]$ mv bad\&filename good_filename
\end{lstlisting}
\end{latin}

در این دستور، قرارگیری بک‌اسلش پیش از علامت \cmd{\&} (که در لینوکس معمولاً برای اجرای دستور در پس‌زمینه به کار می‌رود)، معنای فنی آن را خنثی کرده و به دستور \cmd{mv} اجازه می‌دهد تا آن را صرفاً به‌عنوان بخشی از نام فایل مبدأ شناسایی و پردازش کند.

این تکنیک به‌ویژه زمانی که با فایل‌هایی سر و کار دارید که نام آن‌ها به اشتباه یا توسط سیستم‌عامل‌های دیگر نام‌گذاری شده‌اند، بسیار کارگشا است. به یاد داشته باشید که هر نویسه ویژه باید به‌طور جداگانه با یک بک‌اسلش مهار شود.

\textbf{نکاتی در باب نویسه گریز:}

\begin{itemize}
  \item \textbf{درج نویسه بک‌اسلش:} برای نمایش عینی یک بک‌اسلش در خروجی، باید از خودِ این نویسه جهت گریز دادن استفاده کرد؛ به این معنا که باید دو بک‌اسلش پیاپی (\cmd{\textbackslash\textbackslash}) تایپ شود.
  \item \textbf{رفتار در تک‌کوتیشن:} درون ساختار تک‌کوتیشن (\cmd{' '})، بک‌اسلش هویت عملیاتی خود را از دست داده و صرفاً به عنوان یک کاراکتر متنی ساده و فاقد معنای فنی پردازش می‌شود.
  \item \textbf{گذر از نام‌های مستعار (\en{Alias Bypass}):} یکی دیگر از کاربردهای حرفه‌ای بک‌اسلش، نادیده گرفتن نام‌های مستعار است. به عنوان مثال، اگر دستور \cmd{ls} به صورت \cmd{ls --color=auto} مستعار شده باشد، با درج یک بک‌اسلش پیش از نام دستور (\cmd{\textbackslash ls})، پوسته تعریف مستعار را نادیده گرفته و دستور اصلی را بدون پارامترهای افزوده‌شده اجرا می‌کند.
\end{itemize}

\section{توالی‌های گریز (Escape Sequences)}

نویسه بک‌اسلش (\cmd{\textbackslash}) فراتر از نقش خود در خنثی‌سازی نویسه‌های خاص، به عنوان پیشوند برای بازنمایی کدهای کنترلی (\en{Control Codes}) نیز به کار می‌رود. این کدها که در واقع ۳۲ کاراکتر نخست استاندارد کدگذاری \en{ASCII} را تشکیل می‌دهند، در ابتدا برای ارسال فرامین به تجهیزات سخت‌افزاری نظیر دستگاه‌های تله‌تایپ (\en{Teletype}) طراحی شده بودند.

امروزه برخی از این کدهای کنترلی برای کاربران بسیار آشنا هستند (مانند نویسه تب یا \en{Tab}، پس‌بَر یا \en{Backspace} و خط جدید یا \en{Newline})، در حالی که برخی دیگر نظیر کاراکتر تهی (\en{Null}) یا پایان انتقال (\en{End of Transmission}) در لایه‌های انتزاعی‌تر و تخصصی‌تر سیستم‌عامل مورد استفاده قرار می‌گیرند.

ایده بهره‌گیری از توالی‌های گریز (\en{Escape Sequences}) از زبان برنامه‌نویسی \en{C} وام گرفته شده و امروزه در بسیاری از زبان‌های توسعه و ابزارهای سیستمی، از جمله پوسته‌های لینوکس، به شکلی گسترده پیاده‌سازی شده است. این سازوکار به کاربران و برنامه‌نویسان اجازه می‌دهد تا کدهای کنترلی و کاراکترهای غیرقابل چاپ را در دل رشته‌های متنی گنجانده و مدیریت کنند.

در جدول زیر، برخی از متداول‌ترین توالی‌های گریز و نقش عملیاتی آن‌ها در محیط ترمینال آورده شده است:

\begin{table}[htbp]
\centering
\caption{توالی‌های گریز متداول (\en{Escape Sequences})}
\label{tab:escape-sequences}
\begin{tabularx}{\textwidth}{c l X}
\toprule
\textbf{توالی گریز} & \textbf{مفهوم فنی} & \textbf{شرح عملکرد} \\
\midrule
\cmd{\textbackslash a} & \en{Alert} & سیگنال هشدار که معمولاً منجر به تولید صدای بوق سیستمی (\en{Beep}) می‌شود. \\
\addlinespace
\cmd{\textbackslash b} & \en{Backspace} & نویسه پس‌بَر که مکان‌نما را یک جایگاه به عقب بازمی‌گرداند. \\
\addlinespace
\cmd{\textbackslash n} & \en{Newline} & ایجاد خط جدید؛ در سیستم‌های یونیکسی باعث انتقال مکان‌نما به ابتدای سطر بعدی می‌گردد. \\
\addlinespace
\cmd{\textbackslash r} & \en{Carriage Return} & بازگشت به ابتدای سطر؛ مکان‌نما را بدون تغییر سطر به ابتدای همان خط فعلی می‌برد. \\
\addlinespace
\cmd{\textbackslash t} & \en{Tab} & ایجاد یک فاصله تب افقی (\en{Horizontal Tab}). \\
\bottomrule
\end{tabularx}
\end{table}

هنگام استفاده از دستور \cmd{echo} در خط فرمان، با افزودن گزینه \cmd{-e} (مخفف \en{enable})، قابلیت تفسیر این توالی‌ها فعال می‌شود. بدون این گزینه، عباراتی مانند \cmd{\textbackslash n} صرفاً به صورت متنی چاپ می‌شوند، اما با فعال‌سازی آن، پوسته آن‌ها را به کدهای کنترلی مربوطه ترجمه می‌کند. علاوه بر این روش، در پوسته \en{Bash} می‌توان این توالی‌ها را در قالب ساختار اختصاصی \cmd{\$' '} نیز به کار برد تا مستقیماً توسط خودِ پوسته تفسیر شوند.

با بهره‌گیری از دستور \cmd{sleep} — که ابزاری ساده جهت تعلیق موقت فرآیندهای پوسته برای بازه زمانی مشخص است — می‌توان یک تایمر شمارش معکوس ایجاد کرد. در نمونه زیر، اجرای دستورات به مدت ۱۰ ثانیه متوقف شده و پس از آن، پیامی به همراه هشدار صوتی پخش می‌گردد:

\begin{latin}
\begin{lstlisting}
sleep 10; echo -e "Time's up\a"
\end{lstlisting}
\end{latin}

در این ساختار، نویسه \cmd{;} به عنوان جداکننده دو فرمان عمل می‌کند. پس از انقضای وقفه ۱۰ ثانیه‌ای، دستور \cmd{echo -e} فراخوانی شده و توالی \cmd{\textbackslash a} را به عنوان یک سیگنال صوتی (\en{Alert}) تفسیر و اجرا می‌کند.

یک روش جایگزین و مدرن برای این عملیات، استفاده از نحو اختصاصی \cmd{\$' '} است. این ساختار به صورت بومی در پوسته (\en{In-shell})، تفسیر توالی‌های گریز را فعال کرده و نیاز به گزینه \cmd{-e} در دستور \cmd{echo} را مرتفع می‌سازد:

\begin{latin}
\begin{lstlisting}
sleep 10; echo "Time's up" $'\a'
\end{lstlisting}
\end{latin}

در این وضعیت، دستور \cmd{echo} دو آرگومان مجزا دریافت می‌کند: رشته متنی \cmd{"Time's up"} و خروجی حاصل از بسط \cmd{\$'\textbackslash a'}. هر دو روش یادشده نتیجه عملیاتی یکسانی را رقم می‌زنند و ابزاری کارآمد برای تعبیه هشدارهای شنیداری در اسکریپت‌های اتوماسیون به‌شمار می‌روند.

\section{جمع‌بندی}

با ارتقای سطح تجربه در محیط عملیاتی پوسته، درخواهید یافت که مفاهیم بسط (\en{Expansion}) و کوتیشن (\en{Quoting}) ابزارهایی حیاتی هستند که به صورت مستمر در تعامل با سیستم به کار گرفته می‌شوند. تسلط دقیق بر این دو مبحث، سنگ‌بنای استفاده حرفه‌ای و مؤثر از پوسته محسوب می‌شود؛ چرا که بدون درک عمیق آن‌ها، رفتار پوسته همواره مبهم و پیش‌بینی‌ناپذیر به نظر خواهد رسید و بخش اعظمی از قدرت عملیاتی آن بلااستفاده می‌ماند. به جرات می‌توان اذعان داشت که این مفاهیم، در زمره بنیادی‌ترین دانش‌هایی قرار دارند که هر متخصص لینوکس باید بر آن‌ها چیره شود.

\section{منابع مطالعه تکمیلی}

صفحه راهنمای \cmd{bash} بخش‌های جامعی در مورد بسط (\en{expansion}) و کوتیشن (\en{quoting}) دارد که این مباحث را با بیانی رسمی‌تر پوشش می‌دهند. همچنین، راهنمای مرجع \en{Bash} نیز فصل‌هایی را به این دو موضوع اختصاص داده است:

\begin{latin}
\url{http://www.gnu.org/software/bash/manual/bashref.html}
\end{latin}